\itemhead{习题3.1 · 第15题（教材p.78）}{XJ-B1-C03-HW-P078-Y1137}
\label{XJ-B1-C03-HW-P078-Y1137}
{\footnotesize 湘教版教材p.78 / PDF第85页 / 习题3.1 / 第15题}
\field{原题}{如图，四边形 $ABCD$ 是边长为 4 的正方形，一个动点 $P$ 从点 $A$ 出发，沿着边 $AB\to BC\to CD\to DA$ 运动，返回到点 $A$ 后停止运动，设点 $P$ 走过的路程为 $x$。
\bankfig{C03-HW078-15}{0.36}
\subq{1}将线段 $AP$ 的长度表示成 $x$ 的函数 $f_1(x)$；
\subq{2}设点 $O$ 是正方形 $ABCD$ 的中心，当点 $A,O,P$ 三点不共线时，将 $\triangle AOP$ 的面积表示成 $x$ 的函数 $f_2(x)$。}
\field{参考答案}{（1）$f_1(x)=x$（$0\le x\le4$）、$\sqrt{16+(x-4)^2}$（$4<x\le8$）、$\sqrt{16+(12-x)^2}$（$8<x\le12$）、$16-x$（$12<x\le16$）。（2）$f_2(x)=x$（$0<x\le4$）、$8-x$（$4<x<8$）、$x-8$（$8<x\le12$）、$16-x$（$12<x<16$）；定义域为 $(0,8)\cup(8,16)$，必须排除三点共线的0、8、16。}
\field{核对说明}{定义域核验：三角形面积函数须排除A、O、P共线的路程0、8、16，不能把退化三角形包括进原问题。}
{\small \field{关联知识}{F07 分段函数、绝对值及分段应用（教材帮p.86、91、92、93 / PDF95、100、101、102）；F03 函数求值、解析式与函数建模（教材帮p.80、87、88、89、90、131 / PDF89、96、97、98、99、140）}
\field{方法与易错点}{先列转折点和各段条件，再建立分段式，检查衔接、定义域与值域。 易错：跨段套公式；超额部分按全量收费；绝对值去号漏区间。}
\field{答案依据}{答案册相关题组：PDF第20、21、22页（册内页19、20、21）；按当前题库原题逐项重算，不按旧版题号机械对套。}}
